\paragraph{族的\trModuliSpace。}\label{sub:moduli-spaces-families}
取定一个\trCorneredBoundary{}的\trExactSymplectic\trManifold{}$M$。\trFamily{}$\scr{S}\to\scr{R}$ 的一组\trLagrange\trLabel，是指为诸\trFiber{}取定的、\trLocallyConstant{}的一组\trLagrange\trLabel。与单个\trSurface{}的情形一样，一旦有了\trLabel{}的选取，自然便可把\trExact\trLagrange\trSubmanifold{}的一个对 $(L_{\zeta,0},L_{\zeta,1})$ 与每个 $\zeta\in\Sigma^+$ 相关联，它们是（在任一\trFiber{}中）与 $\zeta$ 相邻的\trBoundaryComponent{}的\trLabel。
设已为 $\scr{S}\to\scr{R}$ 配备了这样的\trLabel，并配备了一组\trStriplikeEnd。又设在每一对 $(L_{\zeta,0},L_{\zeta,1})$ 上都已选定\trFloerDatum{}$(H_\zeta,J_\zeta)$。这时，我们这个\trFamily{}的一个\trPerturbationDatum{}就是一个对 $(K,J)$，由
\begin{equation*}
K\in\Omega^1_{\scr{S}/\scr{R}}(\scr{S},\mathscr{H})
\qquad\text{与}\qquad
J\in C^\infty(\scr{S},\mathscr{J})
\end{equation*}
组成，其在任一\trFiber{}上的限制都满足\cref{sub:floer-perturbation-data}中的诸条件（包括\eqref{eq:perturbation-compatibility}）。
这个记号可能需要一点解释：$K$ 是诸\trFiber{}上 1-形式的一个\trSmooth\trFamily，取值于 $M$ 上\trHamiltonian{}函数所成的空间 $\mathscr{H}$。换个说法：若以 $T\mathscr{S}^{v}\to\mathscr{S}$ 表示沿\trFiber{}的\trTangentBundle，则可以把 $K$ 想成它的\trDualBundle{}经\trPullback{}到 $\mathscr{S}\times M$ 上的一个\trSection（在 $\mathscr{S}\times\partial M$ 附近以及诸\trStriplikeEnd{}上还带有附加条件）。
取定 $\scr{S}\to\scr{R}$ 的一个\trPerturbationDatum。于是，对每一组 $\{y_\zeta\}_{\zeta\in\Sigma}$（其中 $y_\zeta\in\scr{C}(L_{\zeta,0},L_{\zeta,1})$），都可以考察如下的\trModuliSpace
\begin{equation}\label{eq:family-moduli-space}
\scr{M}_{\scr{S}}(\{y_\zeta\})=
\left\{
\begin{aligned}
&(r,u):\ r\in\scr{R},\ \text{而 }u\in C^\infty(\scr{S}_r,M)\ \text{满足 }
\eqref{eq:inhomogeneous-equation}\ \text{与}\ \eqref{eq:moduli-convergence}\\
&\text{就\trPerturbationDatum{}限制到\trFiber{} $\scr{S}_r$ 上而言}
\end{aligned}
\right\}.
\end{equation}
在 $\scr{R}$ 的某个使\trStriplikeEnd{}保持恒定的小开子集 $U\subset\scr{R}$ 上取局部\trTrivialization，便可定义一个（平凡的）\trBanach\trFiberBundle{}$\mathscr{B}_{\scr{S}|U}\to U$，其诸\trFiber{}是来自\cref{sub:linearization}的 $W^{1,p}$ 映射空间 $\mathscr{B}_{\scr{S}_r}$。还有一个\trBanach\trVectorBundle{}$\mathscr{E}_{\scr{S}|U}\to\mathscr{B}_{\scr{S}|U}$，它在每个 $r\in U$ 上的限制是 $L^p$ 空间 $\mathscr{E}_{\scr{S}_r}$。\trInhomogeneous{}的 $\bar\partial$-方程\eqref{eq:inhomogeneous-equation} 给出一个\trSection{}$\mathscr{B}_{\scr{S}|U}\to\mathscr{E}_{\scr{S}|U}$，其零点集就是对一切 $\{y_\zeta\}$ 而言诸 $\mathscr{M}_{\scr{S}|U}(\{y_\zeta\})$ 的并。这个\trSection{}在\trModuliSpace{}一点 $(r,u)$ 处的导数称为\textbf{扩张\trLinearization{}算子}（extended linearized operator）
\begin{equation}\label{eq:extended-linearized-operator}
D_{\scr{S},r,u}\colon (T\mathscr{B}_{\scr{S}|U})_{(r,u)}
=T\scr{R}^{d+1}_r\times(T\mathscr{B}_{\scr{S}_r})_u
\longrightarrow(\mathscr{E}_{\scr{S}_r})_u.
\end{equation}
其中第二个分量就是\eqref{eq:linearization} 的寻常\trLinearization{}算子 $D_{S,u}$（取 $S=\scr{S}_r$），而第一个分量则是\eqref{eq:inhomogeneous-equation} 在保持 $u$ 不动时关于 $r$ 变动的导数。注意，这个扩张后的算子仍是\trFredholm{}的。对 $(r,u)$ 而言，合适的\trRegularity{}概念是 $D_{\scr{S},r,u}$ 的\trSurjective{}性（而不是更强的要求，即 $D_{\scr{S}_r,u}$ 的\trSurjective{}性——后者乃是 $u$ 作为\trDomain{}固定的\trInhomogeneous\trPseudoholomorphicMap{}时的\trRegularity）。照旧，若\trRegularity{}条件对所涉及的一切 $(r,u)$ 都成立，我们就称\trModuliSpace{}$\scr{M}_{\scr{S}}(\{y_\zeta\})$ 或\trPerturbationDatum{}$(K,J)$ 是\trRegular{}的。
\begin{remark}\label{thm:trivialization-independence}
我们的表述绕开了每逢考察\trDomain{}变动的\trPseudoholomorphicMap{}便会出现的一个技术性困难：对整个\trFamily{}$\scr{S}\to\scr{R}$ 而言，不存在\trCanonical{}定义的\trBanachManifold{}$\mathscr{B}_{\scr{S}}$，因为\trTrivialization{}的更换在局部定义的空间 $\mathscr{B}_{\scr{S}|U}$ 之间只诱导\trBanach\trHomeomorphism（而非可微映射）。不过，\trEllipticRegularity{}保证算子\eqref{eq:extended-linearized-operator}——而且在\trRegular{}的情形，还有 $\scr{M}_{\scr{S}}(\{y_\zeta\})$ 的\trSmoothStructure——不依赖于局部\trTrivialization{}的选取。
\end{remark}
